\section{A strict majority on a circle}

\begin{theorem}[982-C11 and the strict convex part of 982-C10]
If \(P\) is a strictly convex planar counterexample and \(m\) of its \(n\) points lie
on a circle, then \(m\le n/2+O(n/\log(2n))\), with an absolute constant.
At every size \(m\ge 2(n-m)+42\) forces a qualifying pin on the circle.
\end{theorem}

\begin{proof}
Set \(k=n-m,\ \Delta=m-k>0\), as there is nothing to prove for the asymptotic statement if \(m\le n/2\).
Normalize the circle to the unit circle and one vertex of its set \(A\) to \(1\).
If \(A\) is contained in a closed semicircle, an end vertex sees \(m-1\) distances along it, reaching the target. Otherwise the center is inside \(\operatorname{conv} A\), and cannot be another vertex.
The group generated by \(A\) has free rank \(r\). The images of \(A\) modulo torsion affinely span \(\mathbb R^r\). For a pin in an extreme fiber of size \(t\), the relation \(bc=a^2\) for equal chord lengths forces both targets of each coincidence to be in the same fiber. The pin therefore sees at least \(m-\lceil t/2\rceil\) radii. If only the circle pins are assumed to fail, every extreme fiber has \(t\ge\Delta+1\), and
\[
 (r+1)(\Delta+1)\le m. \tag{A}
\]
For the asymptotic claim all pins fail. Each off-circle pin \(q\) has a repeated radius among its \(m\) circular targets, so \((q/|q|)^2=bc\) for some \(b,c\in A\). All directions belong to the group \(\Gamma\) of square roots of the group generated by \(A\), still of rank \(r\).

\paragraph{Uniform equalities at different radii.}
Apply Evertse--Schlickewei--Schmidt \cite{ess} to sums of fixed numbers of Laurent monomials with variables in \(\Gamma^d\). Divide each minimal vanishing subsum by one of its monomials. Each resulting nondegenerate equation in ratios has at most \(\exp(O(r+1))\) solutions, since the rank of the monomial image is at most \(dr\). Fixing those ratios cuts out a translate of a finite-component subtorus. Its number of components depends only on the exponent matrix by Smith normal form. Each component with a group point has such a point as a representative. Thus the solutions are covered by \(\exp(O(r+1))\) translates on which the equation is identically zero.
Put \(M=\exp(C(r+1))\), for a fixed \(C\) large enough for the few exponent patterns below. The equation
\[
 \alpha U+\beta U^{-1}+\gamma V+\delta V^{-1}+e=0
\]
with all five coefficients nonzero has no positive-dimensional component of this kind: if both coordinates vary, the constant cannot cancel; if just one varies, its distinct character and inverse cannot cancel. Thus there are at most \(M\) solutions. Squaring distances gives this equation for a fixed noncentral pin and an unknown pair of targets at two fixed unequal radii. It also gives this equation for a fixed target \(q\ne0\), a variable unit-circle pin \(z\) and an unknown target \(Tvz\) on a radius \(T\ne|q|\), using variables \(z,v\). In the latter case at most \(M\) circle pins create any such coincidence.

\paragraph{Double witnesses.}
Fix an off-circle target \(q\) and rotate so \(q=R>0\). If another off-circle target \(b\) has the same distance as \(q\) at two different circle pins \(z,w\), reflection in their chord gives \(b=z+w-Rzw\). With \(\lambda=b/|b|\) set \(X=z/\lambda,\ Y=w/\lambda,\ Z=zw/\lambda\). The reality of \(b/\lambda\) gives
\[
 X-X^{-1}+Y-Y^{-1}=R(Z-Z^{-1}). \tag{B}
\]
On a torus translate where (B) holds, classify the nonconstant coordinate characters up to inversion. Any nonzero class must have at least two coordinates in it. Thus at most one class occurs. A constant triple fixes both pins. If only \(X,Y\) vary, then \(Z=\pm1\) and either \(X+Y=0\) or \(XY=1\); these give respectively equal target radii or collinearity with the center. If only \(X,Z\) or \(Y,Z\) vary, cancellation forces \(R=1\), excluded. If all three vary, write \(Z=st,\ X=xt^{\epsilon_1},\ Y=yt^{\epsilon_2}\) with \(\epsilon_i=\pm1\) and unit modulus coefficients. A positive exponent fixes \(w=Z/X\) or \(z=Z/Y\). When both are negative, comparison of the coefficient of \(t\) gives \(x^{-1}+y^{-1}=-Rs\) and \(xy=s^{-2}\) by conjugating. Substitution gives \(b/q=\sigma-\sigma^2\) with \(\sigma=-s^2t^2\) of unit modulus.

Collect the fixed pins in an exception set \(E_q\) of size at most \(M\), increasing the absolute constant if necessary. Repeating the argument with \(q,b\) interchanged, outside the equal-radius, collinear and exceptional cases both \(v=b/q\) and \(1/v\) would have the last form. Writing \(H=|v|^2\) gives \(2\Re v=3H-H^2=3-1/H\), so \((H-1)^3=0\) and \(b=q\), a contradiction.

\paragraph{Layer counting.}
For \(s\) off-circle targets with \(s_j\) on each radius layer, count losses at circle pins by connecting each distance class of those targets with a simple path. Each unordered pair appears at most twice over all pins. Its first occurrence costs at most \(\binom s2\) in aggregate. Second occurrences in the same radius cost at most \(\sum_j\binom{s_j}{2}\); centered collinear pairs cost at most \(s/2\). Other second occurrences have an exceptional pin on their path edge, and at each \((q,z)\) at most two edges meet \(q\). Thus total loss is at most \(\binom s2+\sum_j\binom{s_j}{2}+(2M+1)s\).

Call a layer large if it contains at least \(\Delta/4\) points; there are \(L\le4n/\Delta\) large layers and \(s\) small-layer points. Count radii including zero. The large layers have at least \((n-s)/2\) distances per pin before merging, as every within-layer multiplicity is at most two. Cross-layer losses at a fixed pin cost at most \(M\binom L2\). Each small-layer target can merge with a large layer at no more than \(M\) circle pins. Since each circle pin sees at most \(n/2\) distances including zero,
\[
 ms/2\le\binom s2+\sum_j\binom{s_j}{2}+(2M+1)s+mM\binom L2+sML .
\]
Here \(m-s\ge\Delta\) and \(s_j<\Delta/4\). The left after moving the first two terms is at least \(3\Delta s/8\). From (A), \(M\), \(s\), and \(L\) are each bounded by \(\exp(O(n/\Delta))\).

For every pin, reflection in its radial axis misses at most \(E=s+2M\binom L2\) targets in any one large layer: every missed point produces a singleton distance before merging the large layers, and the total count is at most \(n/2\). Compose each reflection with that of one fixed pin. There are at least \(n/2\) different rotations since a line through the center contains at most two vertices. On a large layer of size \(v\), each rotation preserves at least \(v-2E\) points. Distinct rotations cannot send the same ordered pair to one another, so \((n/2)(v-2E)\le v^2\). Either \(v\le4E\) or \(v\ge n/2-4E\). If \(\Delta\le8E\), then \(n=(n/\Delta)\Delta\le\exp(O(n/\Delta))\). Otherwise every large layer other than \(A\) falls into the first option, since all off-circle targets together number \(k=(n-\Delta)/2<n/2-4E\). Thus \(k\le s+4EL\le\exp(O(n/\Delta))\).

At a circle pin, start with at least \(\lfloor m/2\rfloor\) circle distances then adjoin the \(k\) targets. At least \((k+1)/2\) counts must merge if the pin fails. Off-circle targets merging with circle distances cost at most \(Mk\) in aggregate by the five-term equation with \(T=1\); off-circle pairs cost at most \(2\binom k2\). This gives \(m(k+1)/2\le Mk+k(k-1)\), hence \(m\le2M+2k\). In all cases \(n\le\exp(O(n/\Delta))\), proving the asymptotic estimate.

For the explicit threshold alone assume just that the circle pins fail, so (A) still applies and \(m\ge2k+42\) forces \(r=0\). We can bound the cost of adding an off-circle target against existing circle distances directly. With \(u=a/z\) for a circle target \(a\) and pin \(z\), the equation is
\[
 u+u^{-1}-\bar q z-qz^{-1}+|q|^2-1=0.
\]
This has finitely many torsion solutions by the same five-term argument applied to the rank-zero group of all roots of unity. Its Newton polygon has area two, so Beukers--Smyth \cite{bs} bounds the number by 44. The solution \(u=1\) is impossible; \(u=-1\) occurs at at most two pins by intersection of the sample circle with the radius-two circle centered at \(q\). The other solutions occur in pairs \(u,u^{-1}\), so at most 23 pins have a merge with this target into the circle distances.
Now \(m(k+1)/2\le23k+k(k-1)=k(k+22)\), contradicted by \(m\ge2k+42\).
\end{proof}
